\section{A finite lattice identity}\label{stencil:section}

Assume that the tensor condition holds. We begin the sufficiency proof by
turning its rational tensor equalities into exact cancellation identities
for finitely many affine functions. These identities will let us construct
two copies of the configuration at different scales whose corresponding
points have the same coordinates in different orders.

Recall that $W=\R\times\R^d$ is viewed over $\Q$, and that
$e_i(v,u)=v+a_i\cdot u$. Let $T\in W\otimes_{\Q}W$ be symmetric, with
\[
 (e_i\otimes e_i)T=0\quad(1\leq i\leq s),
 \qquad G(T)=I_d.
\]
A tensor uses only finitely many vectors. Choose a basis
$w_1,\ldots,w_k$ over $\Q$ for a finite-dimensional subspace containing
its factors, and write
\[
 T=\sum_{a,b=1}^k C_{ab}w_a\otimes w_b,
 \qquad C\in\operatorname{Mat}_k(\Q),\quad C=C^{\mathsf T}.
\]
For $\lambda\in\Q^k$, put
\[
 w(\lambda)=\sum_{j=1}^k\lambda_jw_j.
\]
Write $w_j=(v_j,u_j)$ and let $U=(u_1\ \cdots\ u_k)$ be the real
$d\times k$ matrix whose columns are their gradients. The gradient
identity for $T$ is
\begin{equation}\label{stencil:gram}
 UCU^{\mathsf T}=I_d.
\end{equation}

Set $\Gamma=\Z^k$ and define the evaluation lattices
\[
 \Lambda_i=\{\lambda\in\Gamma:e_i(w(\lambda))=0\},
 \qquad V_i=\operatorname{span}_{\Q}\Lambda_i.
\]
Clearing denominators shows that $V_i$ is the kernel of the rational
linear map $\lambda\mapsto e_i(w(\lambda))$ on $\Q^k$.
Thus $\Q^k/V_i$ embeds into $\R$ over $\Q$. Tensoring this injection
with itself remains injective, so the tensor represented by $C$ has zero
image in $(\Q^k/V_i)\otimes_{\Q}(\Q^k/V_i)$.

The next lemma converts these tensor equalities into a finitely
supported rational function on $\Gamma$. Vanishing of its sums on each
coset will record exact cancellation among equal evaluations.

\begin{lemma}[Finite lattice identity]\label{stencil:finite}
Let $k,s\geq1$, let $\Lambda_1,\ldots,\Lambda_s$ be subgroups of
$\Gamma=\Z^k$, and put $V_i=\operatorname{span}_{\Q}\Lambda_i$.
Let $C\in\operatorname{Mat}_k(\Q)$ be symmetric. Suppose that the
image of
$\sum_{a,b}C_{ab}\varepsilon_a\otimes\varepsilon_b$
in $(\Q^k/V_i)\otimes_{\Q}(\Q^k/V_i)$ is zero for every $i$,
where $\varepsilon_1,\ldots,\varepsilon_k$ is the standard basis of
$\Q^k$. There is a finitely supported function $\nu:\Gamma\to\Q$
such that
\begin{equation}\label{stencil:cosets}
 \sum_{\gamma\in q}\nu(\gamma)=0
 \qquad(q\in\Gamma/\Lambda_i,\ 1\leq i\leq s)
\end{equation}
and
\begin{equation}\label{stencil:moments}
 \sum_\gamma\nu(\gamma)=0,\qquad
 \sum_\gamma\nu(\gamma)\gamma=0,\qquad
 \sum_\gamma\nu(\gamma)\gamma\gamma^{\mathsf T}=C.
\end{equation}
\end{lemma}

\begin{proof}
The Laurent polynomial ring
\[
 R=\Q[X_1^{\pm1},\ldots,X_k^{\pm1}]
\]
is the group algebra of $\Gamma$, with $X^\gamma=\prod_jX_j^{\gamma_j}$.
The quotient map $\Gamma\to\Gamma/\Lambda_i$ induces a ring
homomorphism whose kernel is
\[
 J_i=(X^h-1:h\in\Lambda_i).
\]
Indeed, its kernel consists of Laurent polynomials with zero coefficient
sum on each coset. Within a coset, choose a representative $\gamma_0$
from the finite support; a zero-sum polynomial on that coset is a sum of
terms $c_\gamma(X^\gamma-X^{\gamma_0})$, each in $J_i$.

We will find a Laurent polynomial in every $J_i$ with prescribed
terms through degree two near $X=(1,\ldots,1)$.
To use the tensor hypothesis, we pass to formal logarithmic coordinates,
in which the generators $X^h-1$ become linear forms multiplied by units.
Put $\mathfrak m=(X_1-1,\ldots,X_k-1)$ and complete $R$ with respect
to $\mathfrak m$. With $Z_j=X_j-1$, its completion is
$\widehat R=\Q[[Z_1,\ldots,Z_k]]$: the factors $1+Z_j$ are already
invertible modulo each power of $(Z_1,\ldots,Z_k)$.
The formal changes of variables
\[
 Y_j=\log(1+Z_j),\qquad Z_j=\exp(Y_j)-1
\]
identify this ring with $\Q[[Y_1,\ldots,Y_k]]$ and identify
$\mathfrak m\widehat R$ with $(Y_1,\ldots,Y_k)$.
For each $h\in\Gamma$,
\[
 X^h-1=\exp(h\cdot Y)-1
       =(h\cdot Y)\sum_{n\geq0}\frac{(h\cdot Y)^n}{(n+1)!}.
\]
The final factor is a unit. Hence
\begin{equation}\label{stencil:linear-ideals}
 J_i\widehat R=(h\cdot Y:h\in\Lambda_i)\widehat R.
\end{equation}

Consider the quadratic polynomial
\[
 Q_2(Y)=\tfrac12Y^{\mathsf T}CY.
\]
Identify $\Q[Y_1,\ldots,Y_k]$ with the symmetric algebra of $\Q^k$
by sending $\varepsilon_j$ to $Y_j$. The map to the symmetric algebra
of $\Q^k/V_i$ has kernel generated by $V_i$ in degree one. This follows, for example,
by extending a basis of $V_i$ to a basis of $\Q^k$.
The tensor hypothesis therefore puts $Q_2$ in the ideal generated by
the linear forms $h\cdot Y$, $h\in\Lambda_i$. By
\eqref{stencil:linear-ideals}, $Q_2\in J_i\widehat R$ for every $i$.

The ring $R$ is Noetherian, being a localization of a polynomial ring
over $\Q$. Its $\mathfrak m$-adic completion is therefore flat over
$R$ \cite[Lemma~10.97.2, Tag~00MB]{stacks-completion}.
If $I=\bigcap_{i=1}^sJ_i$, the injection
\[
 R/I\longrightarrow\bigoplus_{i=1}^sR/J_i
\]
remains injective after tensoring with $\widehat R$. Consequently
\[
 I\widehat R=\bigcap_{i=1}^sJ_i\widehat R,
\]
and $Q_2\in I\widehat R$.

We now recover a finite Laurent polynomial while retaining exact
membership in $I$. Write
\[
 Q_2=\sum_{a=1}^N f_a\alpha_a,
 \qquad f_a\in I,\quad\alpha_a\in\widehat R,
\]
as a finite sum. For each $a$, truncate the formal $Z$-expansion of
$\alpha_a$ through degree two to obtain $b_a\in R$ with
$b_a-\alpha_a\in\mathfrak m^3\widehat R$.
Then $f=\sum_af_ab_a\in I$ and
\[
 f(e^{Y_1},\ldots,e^{Y_k})
 =\tfrac12Y^{\mathsf T}CY\pmod{(Y_1,\ldots,Y_k)^3}.
\]
Write the finite Laurent expansion as
$f=\sum_\gamma\nu(\gamma)X^\gamma$.
Membership in each $J_i$ gives \eqref{stencil:cosets}.
Finally,
\[
 X^\gamma=\exp(\gamma\cdot Y)
 =1+\gamma\cdot Y+\tfrac12(\gamma\cdot Y)^2\pmod{(Y)^3}.
\]
Comparing the constant term, gradient, and Hessian at $Y=0$ gives
\eqref{stencil:moments}.
\end{proof}

Apply Lemma~\ref{stencil:finite} to the matrix and evaluation lattices
constructed above. If $\gamma,\gamma'\in\Gamma$, then
$e_i(w(\gamma))=e_i(w(\gamma'))$ exactly when
$\gamma-\gamma'\in\Lambda_i$. Thus \eqref{stencil:cosets} says that
the positive and negative weights cancel separately at every
evaluation value. The second moment in \eqref{stencil:moments}, together
with \eqref{stencil:gram}, gives their prescribed difference in gradient
Gram matrices. We next redistribute these weights so that the negative
part has arbitrarily small gradient second moment, while all these
equalities remain exact.
